\section*{AI工具使用详情}

\subsection*{一、AI工具名称和版本}
本次问题解决过程中使用了ChatGPT（GPT-5.6）、豆包（Seed-1.6-thinking）和
DeepSeek（v3）作为辅助工具。

\subsection*{二、具体使用目的和环节}
\begin{enumerate}[label=（\arabic*）]
  \item \textbf{ChatGPT（GPT-5.6）：}在模型与程序复核环节，用于辅助检查功率平衡、储能SOC递推、
  跨日状态传递、滚动购电结算及不同费用口径；在代码整理环节，用于协助定位测试环境、时间标签
  和实际执行约束中的问题；在论文整理环节，用于LaTeX编译、表格和图表排版检查。
  \item \textbf{豆包（Seed-1.6-thinking）：}在论文撰写和语言检查环节，用于辅助优化中文表达、
  章节标题和表图说明，使文字更简洁、清楚。
  \item \textbf{DeepSeek（v3）：}在模型求解与代码复核环节，用于辅助检查数学符号、约束边界和
  程序结构，并对部分计算流程的效率提出修改建议。
\end{enumerate}

\subsection*{三、关键交互记录}
\subsubsection*{（一）ChatGPT（GPT-5.6）}
\noindent\textbf{Query 1：}给出微网购电、光伏、充放电、弃光与储能电量的目标函数和约束，检查
模型是否物理闭合，尤其检查弃光范围、储能索引和充放电互斥条件。

\noindent\textbf{Output摘要：}建议明确时间步长，补充弃光功率不能超过当期光伏功率的约束，
区分144个功率区间和145个储能状态，并保持充放电二元互斥。

\noindent\textbf{Query 2：}检查问题二实际回测中计划充放电与实际执行是否一致，要求禁止无法
被负荷吸收的放电，并使下一日SOC继承上一日实际日末SOC。

\noindent\textbf{Output摘要：}建议将计划充、放电作为实际动作上限，使用实际充放电递推SOC，
并分别核算计划购电费用、模型期望紧急购电费用和实际紧急购电费用。

\noindent\textbf{Query 3：}依据LaTeX编译日志和逐页预览，检查长表、代码清单、中文字体、图例
遮挡、参考文献与附录分页。

\noindent\textbf{Output摘要：}给出表格换行、代码字体、图例位置和附录拆分建议，并要求通过
重复编译与页面渲染确认结果。

\subsubsection*{（二）豆包（Seed-1.6-thinking）}
\noindent\textbf{Query：}在不改变模型、数值和结论的前提下，检查摘要、章节标题及结果说明的
中文表达，使行文简洁并保持学术风格。

\noindent\textbf{Output摘要：}对部分重复句式、过长句子和表图说明提出精简建议。

\subsubsection*{（三）DeepSeek（v3）}
\noindent\textbf{Query：}检查随机调度与滚动优化代码中的变量定义、约束索引、SOC边界及结算
流程，并分析可以改进的程序结构。

\noindent\textbf{Output摘要：}给出约束组织、变量命名、结果核验和重复计算控制方面的建议。

\subsection*{四、采纳和人工修改情况}
\begin{enumerate}[label=（\arabic*）]
  \item \textbf{ChatGPT（GPT-5.6）：}对于物理约束和实际结算建议，参赛队先重新核对题意、公式
  与现有程序，只采纳能够由约束和审计指标验证的部分；修改后通过功率平衡、SOC递推、SOC上下界、
  跨日连续性、无同时充放电、无虚假放电和费用恒等式检查。对于LaTeX和图表建议，仅调整字体、
  换行、尺寸、图例位置和页面结构，不改变底层数据与结论，并通过完整编译和逐页渲染核验。
  \item \textbf{豆包（Seed-1.6-thinking）：}语言建议由参赛队逐句筛选。对可能改变技术含义、费用
  口径或结论方向的表达不予采纳；保留的修改均再次与公式、表格和结果文件对照。
  \item \textbf{DeepSeek（v3）：}程序结构和效率建议由参赛队结合现有实现人工改写，没有直接
  替换核心模型。涉及求解结果的修改均经过测试、验收程序、年度求解状态和独立审计文件核验。
\end{enumerate}

本参赛队始终主导赛题理解、模型选择、数学推导、参数确定、程序运行、结果判断和论文结论。
所有AI输出均经过人工阅读、修改和核验，未经审查的输出未直接作为核心建模成果提交。
